\section{Python Virtual Environments}

Python includes a built-in mechanism for creating isolated environments.

Virtual environments provide many of the same benefits as Conda and are often simpler for Python-only projects.

\subsection{Creating a Virtual Environment}

Create a new environment called \texttt{venv}:

\begin{lstlisting}
python -m venv venv
\end{lstlisting}

This creates a directory containing an isolated Python installation.

\subsection{Activating the Environment}

Activate the environment:

\begin{lstlisting}
source venv/bin/activate
\end{lstlisting}

Your command prompt should change to indicate that the environment is active.

\subsection{Installing Packages}

Use \texttt{pip} to install packages:

\begin{lstlisting}
pip install numpy
\end{lstlisting}

You may verify the installation:

\begin{lstlisting}
python -c "import numpy"
\end{lstlisting}

\subsection{Leaving the Environment}

Deactivate the environment:

\begin{lstlisting}
deactivate
\end{lstlisting}

\subsection{When Should I Use Conda?}

Conda is often easier when working with:

\begin{itemize}
\item JupyterLab
\item CUDA-enabled Python packages
\item packages with compiled dependencies
\item mixed Python and non-Python software
\end{itemize}

\subsection{When Should I Use a Virtual Environment?}

Virtual environments are often a good choice for:

\begin{itemize}
\item lightweight Python projects
\item software development
\item projects using only pip-installed packages
\end{itemize}

\subsection{Exercise}

Create a virtual environment, install NumPy, verify that it imports successfully, and then deactivate the environment.

\subsection{Recommended Practice}

Create a separate environment for each project.

For example:

\begin{lstlisting}
arctic_project_env
oceananigans_env
machine_learning_env
\end{lstlisting}

Avoid installing large numbers of unrelated packages into a single environment.